\documentclass{article}
\usepackage[
  a4paper,
  margin=0.7in,
  headsep=4pt, % separation between header rule and text
]{geometry}
\usepackage[LGR, T1]{fontenc}
\usepackage{fontspec}
\setmainfont{Times New Roman}
\usepackage{graphicx}
\usepackage{wrapfig}
\usepackage{multicol}
\usepackage{vwcol} 
\usepackage{tikz}
\usepackage{amsmath,bm}
\usepackage[numbers]{natbib}
\usepackage{subcaption}
\usepackage{amsmath,amssymb}
\usepackage{mathtools}
\usepackage[cal=cm]{mathalfa}
\usepackage{geometry}
\usepackage{algorithm}
\usepackage[noend]{algpseudocode}
\usetikzlibrary{arrows,shapes.gates.logic.US,shapes.gates.logic.IEC,calc}
\usetikzlibrary{automata,positioning}
\usepackage{lineno,hyperref}
\usepackage{verbatim}
\usepackage{xcolor}
\usepackage{chngcntr}
\counterwithin{figure}{section}
\usepackage{hyperref}
\usepackage{circuitikz}
\usepackage{pgfplotstable}
\usepackage{pgfplots}
\usepackage{csvsimple}

\newcommand{\HRule}[1]{\rule{\linewidth}{#1}}
\captionsetup[figure]{labelfont={bf},name={Εικόνα},labelsep=period}
\captionsetup[table]{labelfont={bf},name={Πίνακας},labelsep=period}

\title{ \normalsize \textsc{}
		\\ [2.0cm]
		\HRule{0.5pt} \\
		\LARGE \textbf{\uppercase{Eργαστηριακη Αναφορα} \\ [0.5cm]
		\Large Μάθημα : Ηλεκτρονική \\ [0.5cm]
		\large Εργαστήριο : 1ο
		\HRule{2pt} \\ [0.5cm]
		\normalsize \today \vspace*{5\baselineskip}}
		}

\date{}

\author{
		Πάνος Χρηστίδης, ΑΜ:217108 \\
		Ηρακλής Δοξιάδης, ΑΜ:21712 \\
		Αναστάσιος Πεπονής, ΑΜ:21718 \\ [0.5cm]
		Χαροκόπειο Πανεπιστήμιο \\
		Τμήμα Πληροφορικής και Τηλεματικής}

\setlength{\parskip}{1em}

\begin{document}
\maketitle
\thispagestyle{empty}
\newpage

\section{Σκοπός}
\color{blue}
\textit{Στην παρούσα ενότητα περιγράφετε τον σκοπό του εργαστηρίου που παρακολουθήσατε. Η περιγραφή δεν θα πρέπει να ξεπερνάει τις 4-5 γραμμές κειμένου. Προσοχή να δώσετε μία σωστή περιγραφή του εργαστηρίου βασισμένοι και στο φυλλάδιο του εργαστηρίου. Αφαιρέστε το επεξηγηματικό μπλε κείμενο τόσο στην ενότητα αυτή όσο και στις επόμενες πριν παραδώσετε την εργασία}.
\color{black}

\section{Όργανα που χρησιμοποιήθηκαν}
\color{blue}
\textit{Στην παρούσα ενότητα περιγράφετε τα όργανα που χρησιμοποιήσατε συμπεριλαμβάνοντας τόσο τον τύπο όσο και το μοντέλο του οργάνου. Αφαιρέστε το επεξηγηματικό μπλε κείμενο τόσο στην ενότητα αυτή όσο και στις επόμενες πριν παραδώσετε την εργασία}.
\color{black}

Τα όργανα που χρησιμοποιήθηκαν για την παρούσα εργαστηριακή άσκηση ήταν τα εξής:
\begin{itemize}
\item Γεννήτρια συχνοτήτων Tecktronics HFG7746
\item Ψηφιακός παλμογράφος Αgilent DSO9384
\end{itemize}
\section{Κυκλώματα που μελετήθηκαν}
\label{sec:circuits}
\color{blue}
\textit{Στην παρούσα ενότητα περιγράφετε εν συντομία τα κυκλώματα που μελετήσατε στο εργαστήριο στα πλαίσια της εργαστηριακής άσκησης. Αφαιρέστε το επεξηγηματικό μπλε κείμενο τόσο στην ενότητα αυτή όσο και στις επόμενες πριν παραδώσετε την εργασία}.
\color{black}

Τα κυκλώματα που μελετήθηκαν στα πλαίσια της παρούσας εργαστηριακής άσκησης ήταν τα εξής:
\begin{itemize}
\item Γέφυρα διόδων
\item Γέφυρα διόδων με προσθήκη πυκνωτή
\end{itemize}

\section{Κύκλωμα 1}
\color{blue}
\textit{Στην παρούσα ενότητα περιγράφετε αναλυτικά τα κυκλώματα που αναφέρατε στην ενότητα \ref{sec:circuits}. Αλλάξτε τον τίτλο στην παρούσα ενότητα από "Κύκλωμα 1" στο όνομα του κυκλώματος στο οποίο αναφέρεστε (π.χ. Γέφυρα διόδων). Δημιουργήστε όσα sections και αντίστοιχα subsections όπως αυτά που φαίνονται παρακάτω απαιτούνται. Αφαιρέστε το επεξηγηματικό μπλε κείμενο τόσο στην ενότητα αυτή όσο και στις επόμενες πριν παραδώσετε την εργασία}.
\color{black}
\subsection{Περιγραφή κυκλώματος}
\color{blue}
\textit{Στην παρούσα ενότητα περιγράψτε την λειτουργία του κυκλώματος δίνοντας και το σχεδιάγραμμα του χρησιμοποιώντας το περιβάλλον circuitikz του LATEX (\url{https://www.overleaf.com/learn/latex/CircuiTikz_package}). 
Αφαιρέστε το επεξηγηματικό μπλε κείμενο τόσο στην ενότητα αυτή όσο και στις επόμενες πριν παραδώσετε την εργασία}.
\color{black}
\begin{figure}[ht]
	\centering
	\begin{circuitikz} [american]
		\draw
		(0,4) to[R, l=$R$, v=$v_\mathrm{R}$, i=$i_\mathrm{R}$, o-o,color] (0,0) node [below] {D}
		(-2,2) to[D, l=$D_1$, o-o] (0,4) node [above] {B}		
		(0,0) to[D, l=$D_2$, o-o] (2,2) node [above] {C}
		(2,2) to[D, l=$D_3$, o-o] (0,4)
		(0,0) to[D, l=$D_4$, o-o] (-2,2) node [above] {A}
		(2,2) -- (2,-1) -- (-4,-1)
		(-2,2) -- (-4,2) node [above] {A}
		(-4,2) to[sI, v=$v_\mathrm{IN}$, o-o] (-4,-1)  node [below] {C};					
	\end{circuitikz}
	\caption{Σχεδιάγραμμα κυκλώματος προς μελέτη.}
	\label{fig:bridgeD1}
\end{figure}

\subsection{Μετρήσεις που έλαβαν χώρα}
\color{blue}
\textit{Στην παρούσα ενότητα περιγράφετε τις μετρήσεις που λάβατε για το συγκεκριμένο κύκλωμα. Θα πρέπει να δώσετε τόσο τα μεγέθη των μετρήσεων σε μορφή πίνακα και σε γραφική παράσταση (εκτός αν πρόκειται για DC μετρήσεις ή άλλες πιο απλές μετρήσεις. Αφαιρέστε το επεξηγηματικό μπλε κείμενο τόσο στην ενότητα αυτή όσο και στις επόμενες πριν παραδώσετε την εργασία}.
\color{black}
\begin{table}[ht]
\centering
\csvreader[tabular=|c|c|c|,
table head=\hline $V_\mathrm{D}\text{ [mV]}$ & $I_\mathrm{D}\text{ [mA]}$ & $P_\mathrm{D}\text{ [mW]}$ \\\hline\hline,
late after line=\\\hline]
{measurements.csv}{volt=\volt,curr=\curr,power=\power}
{\volt & \curr & \power}
\caption{Μετρήσεις για το κύκλωμα της γέφυρας}
\label{tab:measurements}
\end{table}

\color{blue}
\textit{Αποθηκεύστε τις μετρήσεις σας σε ένα αρχείο .csv και χρησιμοποιήστε την τεχνική του πίνακα \ref{tab:measurements} για να παραθέσετε τις μετρήσεις στην αναφορά. Πάντα να αναφέρετε τις μονάδες.
Αφαιρέστε το επεξηγηματικό μπλε κείμενο τόσο στην ενότητα αυτή όσο και στις επόμενες πριν παραδώσετε την εργασία}.
\color{black}

\begin{figure}
\centering
\begin{tikzpicture}
\begin{axis}[
	xlabel=$V_\mathrm{D}\text{ [mV]}$ , % \hertz requires SIunits
	ylabel=$I_\mathrm{D}\text{ [mA]}$,
	title={Χαρακτηριστική $I_\mathrm{D}=f(V_\mathrm{D})$},
	grid=both,
	minor grid style={gray!25},
	major grid style={gray!25},
	width=0.5\linewidth,
	mark size=2.0pt]
\addplot [line width=1pt,solid,color=blue,mark=*] table [x=volt, y=curr, col sep=comma] {measurements.csv};
\end{axis}
\end{tikzpicture}
\caption{Μετρήσεις για το κύκλωμα της γέφυρας}
\label{fig:bridgeD}
\end{figure}

\color{blue}
\textit{Χρησιμοποιείστε την τεχνική της εικόνας \ref{fig:bridgeD} για να κάνετε τις γραφικές παραστάσεις των μετρήσεων σας. Πάντα να αναφέρετε τις μονάδες.
Αφαιρέστε το επεξηγηματικό μπλε κείμενο τόσο στην ενότητα αυτή όσο και στις επόμενες πριν παραδώσετε την εργασία}.
\color{black}
\section{Συμπεράσματα}
\color{blue}
\textit{Στην παρούσα ενότητα περιγράφετε τα συμπεράσματα των μετρήσεων που κάνατε στα πλαίσια του εργαστηρίου που παρακολουθήσατε. Προσοχή να αναλύσετε προσεκτικά τα συμπεράσματα σας λαμβάνοντας υπόψη και την αντίστοιχη θεωρία καθώς και το φυλλάδιο του εργαστηρίου. Αφαιρέστε το επεξηγηματικό μπλε κείμενο τόσο στην ενότητα αυτή όσο και στις επόμενες πριν παραδώσετε την εργασία}.
\color{black}
\end{document}
